\documentclass[10pt,a4paper]{amsart}
\usepackage[margin=19mm]{geometry}
\usepackage{amsmath,amssymb,mathtools}
\usepackage[scr=rsfs,cal=euler]{mathalfa}
\usepackage{xfrac}
\usepackage[unicode,hidelinks,bookmarks=false]{hyperref}
\newcommand{\ie}{i.e.\ }
\newcommand{\cf}{cf.\ }
\newcommand{\resp}{respectively\ }
\newcommand{\Subsection}{\subsection}
\providecommand{\nobreakdash}{\nobreak\hbox{-}}
\setlength{\emergencystretch}{2em}
\multlinegap0pt
\usepackage{amssymb}
\usepackage{float}
\usepackage{xcolor}
\newtheorem{theorem}{Theorem}
\title{\bf Nearest Graph Laplacians with Prescribed Connected Components: A Convex Framework for Network Reconstruction}
\author{Udit Raj, Sudeepto Bhattacharya, Prince Kanhya}
\date{Source: arXiv:2608.18128v1 (CC BY 4.0)}
\begin{document}
\maketitle

\noindent\emph{Source and licence.} \bf Nearest Graph Laplacians with Prescribed Connected Components: A Convex Framework for Network Reconstruction. Version arXiv:2608.18128v1 (CC BY 4.0), distributed by arXiv under the Creative Commons Attribution 4.0 International licence, http://creativecommons.org/licenses/by/4.0/, as recorded in the arXiv licence metadata for this exact version and shown on the abstract page https://arxiv.org/abs/2608.18128v1. Full licence notice: \texttt{licenses/arxiv-2608.18128-CC-BY-4.0.txt}.\par\smallskip

\noindent\textit{Editorial scope.} Counted unit: the article's theorem theorem (source0.tex lines 710--735). The article's complete original proof of it is delivered with it (source0.tex lines 737--887, one \texttt{proof} environment of 151 lines). No mathematics is written, repaired or re-derived: the statement and the proof are the article's own lines, in the article's own notation, and no retained line is substituted or re-flowed.  No further source-local context is needed: every cross-reference of every retained line resolves inside the two delivered spans.  Nothing was omitted from any retained span, so the package carries no omission record.  No retained line contains a citation, so the wrapper carries no bibliography and no bibliography entry.  No retained line contains a figure environment, an image-inclusion command or drawing code, so no image is delivered and the package declares no asset dependency.  Editorial formatting only, in addition: an AMS \texttt{amsart} preamble that restates the article's own declarations and macros used by the retained lines, namely the environments \texttt{theorem} on separate counters, together with the standard packages those lines need.  The retained environments print in this document as Theorem 1 in order of appearance, whereas the article prints the same items with the numbering of its own full document.  The complete native source of the article as distributed in the arXiv e-print for this exact version is bundled unchanged as \texttt{sources/207970/source0.tex}.
\begin{theorem}[Nearest Laplacian with exactly prescribed components]
	Let $L\in\mathbb R^{n\times n}$ be a given graph Laplacian, and let
	\[
	V=C_1\cup C_2\cup\cdots\cup C_k,
	\qquad
	C_i\cap C_j=\emptyset\quad (i\ne j),
	\]
	be a prescribed partition of the vertex set into nonempty blocks. For each block $C_j$, let $u_j$ be its indicator vector, and define
	\[
	U=[u_1,u_2,\ldots,u_k].
	\]
	Let $M^\star$ be the unique minimizer of problem $(P)$.
	
	Then $M^\star$ is a weighted graph Laplacian whose associated graph has exactly the prescribed connected components
	\[
	C_1,\ldots,C_k.
	\]
	Moreover,
	\[
	\ker(M^\star)=\operatorname{span}\{u_1,u_2,\ldots,u_k\}.
	\]
	If $k<n$, then
	\[
	\lambda_{k+1}(M^\star)\ge \varepsilon>0.
	\]
\end{theorem}

\begin{proof}
	Since $M^\star$ is feasible for $(P)$, it satisfies
	\[
	M^\star U=0.
	\]
	As
	\[
	U=[u_1,u_2,\ldots,u_k],
	\]
	this is equivalent to
	\[
	M^\star u_j=0,\qquad j=1,\ldots,k.
	\]
	Therefore,
	\[
	u_j\in\ker(M^\star),\qquad j=1,\ldots,k.
	\]
	
	Since the prescribed blocks form a partition of the vertex set,
	\[
	\mathbf 1_n=u_1+u_2+\cdots+u_k.
	\]
	Thus,
	\[
	M^\star\mathbf 1_n
	=
	\sum_{j=1}^{k}M^\star u_j
	=
	0.
	\]
	Together with the constraints
	\[
	M^\star=(M^\star)^T,
	\qquad
	M^\star_{ij}\le0\quad (i\ne j),
	\]
	this shows that $M^\star$ is a weighted graph Laplacian.
	
	We next show that there are no edges between distinct prescribed blocks. Fix a block $C_j$ and a vertex $i\in C_j$. Since
	\[
	M^\star u_j=0,
	\]
	the $i$-th component gives
	\[
	\sum_{\ell\in C_j}M^\star_{i\ell}=0.
	\]
	On the other hand, since
	\[
	M^\star\mathbf 1_n=0,
	\]
	we also have
	\[
	\sum_{\ell=1}^{n}M^\star_{i\ell}=0.
	\]
	Subtracting these two equalities yields
	\[
	\sum_{\ell\notin C_j}M^\star_{i\ell}=0.
	\]
	For every $\ell\notin C_j$, we have $i\ne \ell$, and hence
	\[
	M^\star_{i\ell}\le0.
	\]
	A sum of nonpositive numbers can be zero only if every term is zero. Therefore,
	\[
	M^\star_{i\ell}=0,
	\qquad
	i\in C_j,\quad \ell\notin C_j.
	\]
	Thus, there are no edges between distinct prescribed blocks.
	
	Consequently, after a suitable permutation of the vertices, $M^\star$ becomes block diagonal:
	\[
	Q M^\star Q^T
	=
	\operatorname{diag}(M^\star_1,M^\star_2,\ldots,M^\star_k),
	\]
	where
	\[
	M^\star_j=M^\star[C_j,C_j].
	\]
	
	We now prove that each prescribed block is internally connected. Since there are no entries connecting $C_j$ with the other blocks and
	\[
	M^\star u_j=0,
	\]
	the corresponding principal block satisfies
	\[
	M^\star_j\mathbf 1_{n_j}=0.
	\]
	Thus, $\mathbf 1_{n_j}$ is always a zero direction of $M^\star_j$.
	
	From the constraint in $(P)$, we have
	\[
	M^\star_j\succeq \varepsilon P_j.
	\]
	Therefore, for every $x\in\mathbb R^{n_j}$,
	\[
	x^TM^\star_jx\ge \varepsilon x^TP_jx.
	\]
	If $x\perp \mathbf 1_{n_j}$, then $P_jx=x$, and hence
	\[
	x^TM^\star_jx\ge \varepsilon\|x\|_2^2.
	\]
	Thus, for every non-singleton block $C_j$, the matrix $M^\star_j$ is positive definite on
	\[
	\mathbf 1_{n_j}^{\perp}.
	\]
	It follows that
	\[
	\ker(M^\star_j)=\operatorname{span}\{\mathbf 1_{n_j}\},
	\]
	and hence
	\[
	\lambda_2(M^\star_j)\ge\varepsilon>0
	\]
	for every block with $n_j\ge2$. By the standard Laplacian nullity theorem, the subgraph induced by each non-singleton block $C_j$ is connected. Singleton blocks are connected by convention.
	
	Thus, there are no edges between distinct prescribed blocks, and each prescribed block is internally connected. Hence the connected components of the graph associated with $M^\star$ are exactly
	\[
	C_1,C_2,\ldots,C_k.
	\]
	
	Therefore, the graph has exactly $k$ connected components. By the Laplacian nullity theorem,
	\[
	\dim\ker(M^\star)=k.
	\]
	Since
	\[
	u_1,u_2,\ldots,u_k
	\]
	are linearly independent and belong to $\ker(M^\star)$, we conclude that
	\[
	\ker(M^\star)
	=
	\operatorname{span}\{u_1,u_2,\ldots,u_k\}.
	\]
	
	Finally, since $M^\star$ is permutation-similar to
	\[
	\operatorname{diag}(M^\star_1,M^\star_2,\ldots,M^\star_k),
	\]
	its eigenvalues are the union of the eigenvalues of the block matrices $M^\star_j$. Each block contributes exactly one zero eigenvalue. Hence the first $k$ eigenvalues of $M^\star$ are zero. If $k<n$, then at least one block is non-singleton, and
	\[
	\lambda_{k+1}(M^\star)
	=
	\min_{\{j:n_j\ge2\}}\lambda_2(M^\star_j)
	\ge
	\varepsilon.
	\]
	This completes the proof.
\end{proof}

\end{document}
